% Chapter 4: Terms and Definitions
\chapter{术语与定义对比}

\section{术语体系概览}

\subsection{术语数量对比}

\begin{table}[H]
\centering
\caption{三版本术语数量对比}
\begin{tabular}{|l|c|c|c|}
\hline
\textbf{术语类别} & \textbf{\versionAShort} & \textbf{\versionBShort} & \textbf{\versionCShort} \\
\hline
基础术语         & 约 30 条     & 约 35 条     & 约 40 条     \\
\hline
装置类术语       & 约 25 条     & 约 32 条     & 约 38 条     \\
\hline
技术参数术语     & 约 20 条     & 约 28 条     & 约 35 条     \\
\hline
测试相关术语     & 约 15 条     & 约 20 条     & 约 25 条     \\
\hline
智能照明术语     & 0 条        & 约 12 条     & 约 15 条     \\
\hline
\textbf{合计}    & \textbf{约 90 条} & \textbf{约 127 条} & \textbf{约 153 条} \\
\hline
\end{tabular}
\end{table}

\section{关键术语对照}

\subsection{基础几何术语}

\begin{table}[H]
\centering
\caption{基础几何术语对照}
\small
\begin{tabularx}{\textwidth}{|l|X|X|X|}
\hline
\textbf{术语} & \textbf{\versionAShort} & \textbf{\versionBShort} & \textbf{\versionCShort} \\
\hline
透光平面 & 透光平面定义基本与欧版一致 & 同左，新增"对矩阵式 LED 的特殊规定" & Light-emitting surface \\
\hline
参考轴 & 灯具参考中心发出的水平轴 & 同左 & Reference axis, defined precisely \\
\hline
几何可见度 & 角锥可见度测试 & 强化角度定义 & Geometric visibility \\
\hline
对地基准面 & 静止时的水平面 & 同左，新增"对于可变高度车辆的考虑" & Ground reference plane \\
\hline
车辆中心平面 & 车辆纵向对称平面 & 同左 & Median longitudinal plane \\
\hline
\end{tabularx}
\end{table}

\subsection{装置类术语对照}

\begin{table}[H]
\centering
\caption{装置类术语对照}
\small
\begin{tabularx}{\textwidth}{|l|X|X|X|}
\hline
\textbf{术语} & \textbf{\versionAShort} & \textbf{\versionBShort} & \textbf{\versionCShort} \\
\hline
远光灯 & 照射前方远距离道路的灯具 & 同左，新增"ADB 中的远光功能"说明 & Driving beam headlamp \\
\hline
近光灯 & 照射前方近距离道路、不眩目对面车辆的灯具 & 同左 & Passing beam headlamp \\
\hline
DRL & 昼间行驶时提高车辆可见性的灯具 & 强化定义，新增"DRL 与位置灯的协同" & Daytime running lamp \\
\hline
AFS & 根据行驶条件自动调整光束的前照灯系统 & 强化定义，新增"AFS 等级转换要求" & Adaptive front lighting system \\
\hline
ADB & 未单独定义 & 新增：根据前方车辆自动调整远光光型的系统 & Adaptive driving beam \\
\hline
ESS & 在紧急制动时高频率闪烁的制动灯信号 & 强化定义，明确闪烁频率、持续时间等 & Emergency stop signal \\
\hline
矩阵式 LED & 未定义 & 新增：由多个独立可控制 LED 单元组成的前照灯 & Matrix LED \\
\hline
\end{tabularx}
\end{table}

\subsection{光度学术语}

\begin{table}[H]
\centering
\caption{光度学术语对照}
\small
\begin{tabularx}{\textwidth}{|l|X|X|X|}
\hline
\textbf{术语} & \textbf{\versionAShort} & \textbf{\versionBShort} & \textbf{\versionCShort} \\
\hline
发光强度 & 单位立体角内发出的光通量，单位 cd & 同左 & Luminous intensity, in candela \\
\hline
光通量 & 单位时间内发射、传输或接收的光能 & 同左 & Luminous flux, in lumen \\
\hline
照度 & 单位面积上接收的光通量，单位 lx & 同左 & Illuminance, in lux \\
\hline
色温 & 与黑体辐射比较得出的颜色度量 & 同左 & Colour temperature, in Kelvin \\
\hline
色度坐标 & CIE 1931 色度图上的坐标 & 强化定义，新增"CIE 1976 替代选项" & Chromaticity coordinates \\
\hline
基准点 & 灯具上的特定测试参考位置 & 同左 & Reference point \\
\hline
光通维持率 & 老化后光通量与初始光通量的比值 & 强化定义，新增"针对 LED 的快速老化方法" & Lumen maintenance \\
\hline
\end{tabularx}
\end{table}

\subsection{测试相关术语}

\begin{table}[H]
\centering
\caption{测试相关术语对照}
\small
\begin{tabularx}{\textwidth}{|l|X|X|X|}
\hline
\textbf{术语} & \textbf{\versionAShort} & \textbf{\versionBShort} & \textbf{\versionCShort} \\
\hline
基准中心 & 灯具的特征中心点 & 同左 & Centre of reference \\
\hline
光通维持率测试 & 1000 小时老化后测试 & 新增 500 小时加速老化方法 & Lumen maintenance test \\
\hline
色度测试角度 & 多角度（通常为 HV、50L、50R 等） & 同左，新增"ADB 色度测试角度" & Multi-angle chromaticity test \\
\hline
光型 & 灯具投射光线的空间分布 & 同左，新增"矩阵式 LED 光型定义" & Beam pattern definition \\
\hline
眩目 & 对向驾驶员感到刺眼的现象 & 同左 & Glare \\
\hline
测试屏幕 & 距离灯具一定距离的接收屏 & 同左 & Test screen \\
\hline
\end{tabularx}
\end{table}

\section{新增术语详解}

\subsection{ADB 相关术语（草案版新增）}

\begin{description}[leftmargin=2em, style=nextline]
    \item[自适应远光灯（ADB）] 能根据前方车辆、交通环境、道路条件自动调整远光光型的前照灯系统。
    \item[ADB 控制单元] 实现 ADB 功能的电子控制单元，通常包含图像传感器、信号处理模块、驱动模块。
    \item[ADB 光型段（Segment）] ADB 系统中的一个独立可控区域，对应若干颗 LED。
    \item[ADB 切换时间] ADB 系统从识别到目标到完成光型调整所需的时间。
    \item[ADB 识别距离] ADB 系统能可靠识别前方车辆的最大距离。
\end{description}

\subsection{矩阵式 LED 相关术语}

\begin{description}[leftmargin=2em, style=nextline]
    \item[矩阵式 LED 模块（Matrix LED module）] 由多个独立可控的 LED 单元组成的灯具模块。
    \item[像素（Pixel）] 矩阵式 LED 中的单个可独立控制的 LED 单元。
    \item[像素分辨率] 矩阵式 LED 模块在水平和垂直方向上的像素数量。
    \item[失效模式] 矩阵式 LED 中部分像素失效时系统的应对方式。
    \item[降级模式] 部分功能失效时系统仍能维持的基本功能。
\end{description}

\subsection{Logo 灯与迎宾灯相关术语}

\begin{description}[leftmargin=2em, style=nextline]
    \item[Logo 投影灯] 在地面上投射厂商或车型标识的灯具。
    \item[迎宾灯] 车门打开时照亮地面的灯具。
    \item[装饰性照明] 不具备安全功能、仅用于美化的照明。
    \item[图案投影距离] Logo 投影灯投射图案的距离。
\end{description}

\subsection{EMC 相关术语}

\begin{description}[leftmargin=2em, style=nextline]
    \item[电磁兼容性（EMC）] 设备在电磁环境中能正常工作且不对其他设备产生不可接受的电磁干扰的能力。
    \item[辐射发射] 设备通过电磁辐射向外发射的电磁能量。
    \item[传导发射] 设备通过电源线或信号线向外传导的电磁能量。
    \item[抗扰度] 设备抵抗外部电磁干扰的能力。
\end{description}

\section{删除术语分析}

\subsection{现行版特有术语的处理}

现行版中有部分术语在草案版中被删除或整合：

\begin{table}[H]
\centering
\caption{现行版特有术语在草案版中的处理}
\begin{tabular}{|l|l|}
\hline
\textbf{术语} & \textbf{草案版处理} \\
\hline
HID 灯专属术语     & 整合到通用术语，删除 HID 专属条款 \\
\hline
卤素灯专属术语     & 整合，删除独立条目 \\
\hline
"特殊用途车辆" 术语  & 保留并扩展，新增"低速电动车" \\
\hline
"过渡条款" 术语     & 整合到适用范围 \\
\hline
\end{tabular}
\end{table}

\section{术语差异带来的影响}

\subsection{对认证的影响}

术语差异对认证工作产生以下影响：

\begin{enumerate}[leftmargin=2em]
    \item 申请企业需要更新产品技术文件中的术语表述；
    \item 检测报告需要按新术语重新填写；
    \item 认证机构需要发布术语对照指引；
    \item 国际认证对接时需要进行术语映射。
\end{enumerate}

\subsection{对技术研究的影响}

术语统一有利于：

\begin{enumerate}[leftmargin=2em]
    \item 学术研究和论文写作；
    \item 国际技术交流；
    \item 工程师之间的沟通；
    \item 行业报告的标准化。
\end{enumerate}

\section{小结}

\begin{enumerate}[leftmargin=2em]
    \item 草案版术语数量较现行版增加约 40\%，新增术语主要集中在智能照明、EMC、Logo 灯等领域；
    \item 草案版术语定义与国际标准的对齐度从约 90\% 提升至约 95\%；
    \item 草案版保留了中国特色的术语（如"特殊用途车辆"），体现了国情的考量；
    \item 术语的统一和扩展有助于提高标准的可操作性和国际兼容性。
\end{enumerate}